% ============================================================================
% Municipality-Scale Flood Risk Mapping in La Guajira, Colombia
% Using Sentinel-1 SAR and Ensemble Machine Learning (2015--2025)
% Authors: Cristian Espinal Maya, Santiago Jimenez Londono
% Format: arXiv preprint
% Replicated from Antioquia framework — adapted for arid/semi-arid department
% ============================================================================

\documentclass{article}

\usepackage{arxiv}

\usepackage[utf8]{inputenc}
\usepackage[T1]{fontenc}
\usepackage{hyperref}
\usepackage{url}
\usepackage{booktabs}
\usepackage{amsfonts}
\usepackage{amssymb}
\usepackage{amsmath}
\usepackage{nicefrac}
\usepackage{microtype}
\usepackage{graphicx}
\usepackage{multirow}
\usepackage{siunitx}
\sisetup{
    group-separator = {,},
    group-minimum-digits = 4,
}
\usepackage{natbib}
\usepackage{tabularx}
\usepackage{textcomp}

\graphicspath{{./figures/}}

\title{Municipality-Scale Flood Risk Mapping in La Guajira, Colombia, Using Sentinel-1 SAR and Ensemble Machine Learning in an Arid--Semi-Arid Environment (2015--2025)}

\author{
  Cristian Espinal Maya\thanks{Correspondence: cespinalm@eafit.edu.co. ORCID: \href{https://orcid.org/0009-0000-1009-8388}{0009-0000-1009-8388}} \\
  School of Applied Sciences and Engineering\\
  Universidad EAFIT\\
  Medell\'in, Colombia \\
  \texttt{cespinalm@eafit.edu.co} \\
  \And
  Santiago Jim\'enez Londo\~no\thanks{ORCID: \href{https://orcid.org/0009-0007-9862-7133}{0009-0007-9862-7133}} \\
  School of Applied Sciences and Engineering\\
  Universidad EAFIT\\
  Medell\'in, Colombia \\
  \texttt{sjimenez@eafit.edu.co} \\
}

\begin{document}
\maketitle

\begin{abstract}
Flood risk in arid and semi-arid regions is systematically underestimated: global SAR-based flood monitoring products routinely exclude these areas due to sand--water backscatter confusion, and the scarcity of flood events reduces the historical record available for hazard modelling. We address this gap with a municipality-scale flood risk framework for the Department of La Guajira, Colombia (\SI{20848}{\kilo\metre\squared}; $\sim$1.0 million inhabitants)---the country's most arid department and one of its most climate-vulnerable. We processed Sentinel-1 C-band SAR scenes (2015--2025) within Google Earth Engine, implementing a Sand Exclusion Layer (SEL) based on multi-temporal backscatter statistics to resolve the sand--water confusion that has limited SAR flood mapping in desert environments \citep{Martinis2018, Garg2024}. Eighteen predictor variables were integrated into a weighted ensemble of Random Forest, XGBoost, and LightGBM under spatial five-fold cross-validation. The framework quantifies the drought--flood paradox characteristic of La Guajira: La Ni\~na years amplify flood extent relative to El Ni\~no years, while desiccated soils from preceding droughts reduce infiltration capacity and intensify flash-flood response. Overlaying the susceptibility surface with \SI{100}{\metre} population data reveals disproportionate exposure among the Wayuu indigenous community ($\sim$45\% of residents), who occupy the most climate-vulnerable territories. The entire framework---data, code, and trained models---is open-access and fully reproducible, extending the methodology previously developed for the Department of Antioquia \citep{EspinalMaya2026} to demonstrate its applicability in markedly different hydroclimatic environments.
\end{abstract}

\keywords{flood susceptibility mapping \and Sentinel-1 SAR \and ensemble machine learning \and arid regions \and sand exclusion layer \and HAND \and Google Earth Engine \and SHAP \and ENSO \and drought--flood paradox \and Wayuu \and La Guajira Colombia}

%%%%%%%%%%%%%%%%%%%%%%%%%%%%%%%%%%%%%%%%%%
\section{Introduction}
\label{sec:introduction}

In November 2024, Hurricane Rafael struck the Colombian Caribbean coast, triggering floods that affected over 195,000 people across the Department of La Guajira. Municipal authorities lacked spatially explicit flood hazard maps to guide emergency response, relying instead on national-scale products that could not distinguish which arroyos, river reaches, or coastal stretches had been inundated. This event followed a near-identical pattern from October 2022, when Hurricane Julia caused flooding severe enough to declare public calamity in eight of La Guajira's 15 municipalities. Yet between these catastrophic flood events, the same communities endure some of Colombia's most severe droughts---a paradox that makes La Guajira a uniquely challenging and scientifically important case for flood risk assessment.

Globally, flooding affects approximately 1.65 billion people per decade \citep{UNDRR2020}, and \citet{Tellman2021} showed that population exposure to flooding grew 20--24\% between 2000 and 2018. However, flood risk in arid and semi-arid regions remains systematically understudied \citep{Alabbad2023}. The Copernicus Emergency Management Service Global Flood Monitoring system explicitly excludes arid areas from its analyses because C-band SAR backscatter from dry sand surfaces is spectrally indistinguishable from calm water \citep{Garg2024}. This exclusion creates a blind spot precisely where flash floods are most lethal: the combination of low infiltration capacity, sparse vegetation, and high-intensity rainfall produces rapid runoff and extreme flood velocities \citep{Alabbad2023}.

Three developments converge to address this gap. First, \citet{Martinis2018} demonstrated that a Sand Exclusion Layer (SEL) derived from multi-temporal Sentinel-1 statistics can resolve the sand--water confusion, and \citet{Garg2024} showed that fusing VV coherence with amplitude information improves arid flood detection by $\sim$50\%. Second, ensemble machine learning (Random Forest, XGBoost, LightGBM) has been validated for flood susceptibility in semi-arid environments \citep{Bammou2024, Abedi2022}, with SHAP enabling interpretability \citep{Lundberg2017}. Third, Google Earth Engine democratizes petabyte-scale satellite processing \citep{Gorelick2017}, enabling department-scale analyses that would be computationally prohibitive on local infrastructure.

Despite these advances, no study has produced municipality-level flood susceptibility maps for a Colombian arid department. The broader literature on flood susceptibility in tropical arid/semi-arid regions remains sparse: \citet{Bammou2024} mapped flood susceptibility in the Moroccan High Atlas (semi-arid), and \citet{Panahi2022} applied SAR-ML hybrid models to Iran's Karoon River floodplain, but neither addressed the compound drought--flood dynamics that characterize Caribbean arid zones under ENSO influence. For La Guajira specifically, existing flood studies are limited to the city of Riohacha \citep{Nardini2016, PerezMontiel2022}, and no department-wide assessment exists.

\textbf{Research gap and contributions.} This study makes four specific contributions: (1) the first municipality-scale flood susceptibility map for an arid Colombian department, extending the framework of \citet{EspinalMaya2026} to a hydroclimatically contrasting environment; (2) integration of a Sand Exclusion Layer and coastal/tidal masking to overcome SAR limitations in arid and coastal terrain; (3) quantification of the drought--flood paradox through ENSO-stratified analysis of flood extent and the role of antecedent soil moisture; and (4) explicit assessment of Wayuu indigenous community exposure, linking physical hazard mapping with vulnerability dimensions identified by \citet{Contreras2020} and \citet{FernandezLopera2024}.

\subsection{Study Area}
\label{sec:study_area}

\begin{figure}[htbp]
\centering
\includegraphics[width=0.75\textwidth]{fig01_study_area.pdf}
\caption{Department of La Guajira, Colombia, showing three subregions and 15 municipalities. Inset: location within Colombia and proximity to the Caribbean Sea and Venezuela. The R\'io Rancher\'ia corridor (dashed blue line) traverses the department from the Sierra Nevada de Santa Marta to the Caribbean coast.}
\label{fig:study_area}
\end{figure}

The Department of La Guajira (\SI{20848}{\kilo\metre\squared}; 10\textdegree{}22'--12\textdegree{}28'\,N, 71\textdegree{}06'--73\textdegree{}39'\,W) occupies the northernmost peninsula of South America, bounded by the Caribbean Sea (north and west), Venezuela (east), and the departments of Cesar and Magdalena (south) (Figure~\ref{fig:study_area}). With $\sim$1.0 million inhabitants across 15 municipalities grouped into three subregions, the department spans an extraordinary climatic gradient within a small area.

\textbf{Alta Guajira} (Uribia, Manaure; $\sim$\SI{9000}{\kilo\metre\squared}) is a hyper-arid desert (K\"oppen BWh) with annual rainfall of 125--400~mm, extensive sand dune fields, the isolated Serran\'ia de Macuira ($\sim$\SI{865}{\metre}), and Colombia's largest marine salt complex at Manaure. It is the ancestral territory of the Wayuu people, Colombia's largest indigenous group.

\textbf{Media Guajira} (Riohacha, Maicao, Dibulla; $\sim$\SI{5000}{\kilo\metre\squared}) transitions from semi-arid (BSh) to humid coastal plains near Dibulla, with rainfall of 400--800~mm and the R\'io Rancher\'ia delta. The department capital Riohacha sits at the river mouth, where both fluvial and coastal flooding converge.

\textbf{Baja Guajira} (10 municipalities; $\sim$\SI{7000}{\kilo\metre\squared}) encompasses the foothills of the Sierra Nevada de Santa Marta (west) and the Serran\'ia del Perij\'a (east), with tropical savanna climate (Aw), rainfall of 800--2,500+~mm, and the upper reaches of the R\'io Rancher\'ia and R\'io Cesar. This subregion hosts the Cerrej\'on coal mining complex and most of the department's agricultural activity.

The R\'io Rancher\'ia ($\sim$\SI{248}{\kilo\metre}) is the principal hydrological axis, originating on the Sierra Nevada at $\sim$\SI{3875}{\metre} and discharging into the Caribbean near Riohacha. Its torrential regime produces extreme seasonal variability: base flows approach zero during dry months (December--March), while discharge rises dramatically during the September--November wet season, particularly when amplified by La Ni\~na \citep{PerezMontiel2022}. Secondary rivers---the Palomino, Jerez, and Ancho, descending steeply from the Sierra Nevada---produce flash floods in Dibulla. In the east, the R\'io Carraipia flows toward Venezuela.

Precipitation is predominantly unimodal, with the main wet season concentrated in September--November and a prolonged dry season from December through March. This contrasts with the bimodal (MAM/SON) pattern of Andean Colombia. Rainfall is strongly modulated by ENSO: El Ni\~no reduces precipitation by 30--50\% across the department, producing severe droughts \citep{MeraGomez2021}, while La Ni\~na increases precipitation by 20--40\%, triggering floods \citep{Hoyos2013, SedanoCruz2017}. The 2012--2016 El Ni\~no-associated drought killed at least 4,770 Wayuu children under five \citep{Contreras2020}, underscoring the extreme climate vulnerability of the region.

%%%%%%%%%%%%%%%%%%%%%%%%%%%%%%%%%%%%%%%%%%
\section{Materials and Methods}
\label{sec:methods}

\subsection{Data Sources}

The framework integrates 10 satellite and ancillary datasets (Table~\ref{tab:data_sources}), all accessed via Google Earth Engine \citep{Gorelick2017} except administrative boundaries (FAO GAUL).

\begin{table}[htbp]
\caption{Datasets used in the flood risk assessment framework.}
\label{tab:data_sources}
\centering
\footnotesize
\begin{tabular}{lllll}
\toprule
\textbf{Dataset} & \textbf{Source} & \textbf{Resolution} & \textbf{Period} & \textbf{Role} \\
\midrule
Sentinel-1 GRD & ESA/Copernicus & \SI{10}{\metre} & 2015--2025 & Flood detection \\
JRC GSW v1.4 & EC/JRC & \SI{30}{\metre} & 1984--2021 & Water dynamics \\
SRTM DEM v3 & NASA/USGS & \SI{30}{\metre} & 2000 & Topographic features \\
MERIT Hydro v1.0.1 & Yamazaki et al. & \SI{90}{\metre} & 2019 & HAND \\
CHIRPS Daily v2.0 & UCSB/CHG & \SI{5.5}{\kilo\metre} & 2015--2025 & Precipitation \\
ERA5-Land Monthly & ECMWF/C3S & \SI{11}{\kilo\metre} & 2015--2025 & Soil moisture \\
Sentinel-2 MSI & ESA/Copernicus & \SI{10}{\metre} & 2020--2023 & NDVI \\
ESA WorldCover v200 & ESA & \SI{10}{\metre} & 2021 & Land cover \\
WorldPop v2020 & U. Southampton & \SI{100}{\metre} & 2020 & Population density \\
FAO GAUL 2015 & FAO & Vector & 2015 & Admin. boundaries \\
\bottomrule
\end{tabular}
\end{table}

Sentinel-1 C-band SAR GRD scenes were ingested (IW mode, VV+VH polarization, descending orbit). The descending pass (morning acquisition, $\sim$06:00 local time) was preferred for La Guajira because calmer wind conditions reduce wave-induced backscatter on coastal waters and inland flood surfaces, improving water detection accuracy. The Sentinel-1B failure (December 2021 to April 2024) reduced revisit to 12 days; Sentinel-1C deployment in December 2024 restored 6-day coverage. The JRC Global Surface Water dataset \citep{Pekel2016} provided 38-year water dynamics for validation and as an input feature. Terrain variables (elevation, slope, aspect, curvature, TWI, SPI) were derived from SRTM \SI{30}{\metre}. HAND was computed from MERIT Hydro \SI{90}{\metre} \citep{Yamazaki2019} and resampled to \SI{30}{\metre} via bilinear interpolation. While resampling does not add spatial information beyond the native \SI{90}{\metre}, it enables pixel-aligned stacking; the effective resolution remains \SI{90}{\metre}, and its vertical accuracy ($\pm$\SI{5}{\metre}, inherited from SRTM) should be considered when interpreting threshold-based results.

\subsection{SAR-Based Water Detection with Sand Exclusion}
\label{sec:sar_water}

\begin{figure}[htbp]
\centering
\includegraphics[width=0.85\textwidth]{fig02_sar_water_detection.pdf}
\caption{Sentinel-1 SAR flood detection in the R\'io Rancher\'ia corridor, La Guajira. (\textbf{a})~Pre-flood VV backscatter (dry season). (\textbf{b})~During-flood VV (Oct--Nov La Ni\~na event). (\textbf{c})~Sand Exclusion Layer (orange) overlaid on detected flood water (blue). Note the critical role of the SEL in preventing false flood detection over desert sand surfaces that share low backscatter signatures with water.}
\label{fig:methodology}
\end{figure}

The SAR workflow transforms Sentinel-1 backscatter into monthly binary water maps through five steps: (1) speckle reduction via focal median filter (\SI{50}{\metre} kernel) with slope masking ($> 30^{\circ}$); (2) \textbf{Sand Exclusion Layer (SEL)} generation using multi-temporal backscatter statistics; (3) adaptive Otsu thresholding \citep{Otsu1979} on the VV histogram; (4) water mask refinement with \SI{2}{\hectare} minimum mapping unit, permanent water flagging (JRC occurrence $\geq 75\%$), and coastal/tidal masking; and (5) monthly maximum-extent compositing.

\textbf{Sand Exclusion Layer.} The SEL is critical for arid La Guajira, where dry sand and calm water both produce low VV backscatter ($< -15$~dB). Following \citet{Martinis2018}, we computed the temporal minimum and standard deviation of VV backscatter across dry-season scenes (December--March) over the full 2015--2025 record. Pixels with mean dry-season backscatter $< -18$~dB in $\geq$80\% of dry-season composites were classified as persistent sand and excluded from water detection. This approach leverages the temporal stability of sand surfaces versus the transient nature of flood water \citep{Garg2024}.

\textbf{Coastal/tidal masking.} Tidal flats near the Rancher\'ia delta and the Manaure salt pans produce low-backscatter signatures at low tide that mimic flood water. We generated a permanent water/salt-flat mask from the minimum backscatter composite (entire record) and applied a \SI{500}{\metre} coastal buffer in Manaure and Uribia municipalities.

\subsection{Flood Frequency Mapping}

Flood frequency per pixel was computed as:
\begin{equation}
    FF_i = \frac{\sum_{m=1}^{M} W_{i,m}}{M} \times 100
    \label{eq:flood_freq}
\end{equation}
where $M$ is the total number of monthly composites and $W_{i,m}$ is binary water detection. Given La Guajira's lower flood frequency, classes were adjusted: permanent ($>$75\%), very frequent (50--75\%), frequent (15--50\%), occasional (5--15\%), and rare (1--5\%). Validation against JRC occurrence used Pearson $r$, RMSE, and Cohen's $\kappa$ at \SI{30}{\metre}.

\subsection{Flood Susceptibility Modeling}

Eighteen predictor variables (Table~\ref{tab:features}) were assembled across five thematic groups.

\begin{table}[htbp]
\caption{Predictor variables organized by thematic group.}
\label{tab:features}
\centering
\footnotesize
\begin{tabular}{lll}
\toprule
\textbf{Group} & \textbf{Variable} & \textbf{Source} \\
\midrule
\multirow{7}{*}{Topographic} & Elevation, Slope, Aspect & SRTM \\
 & Plan curvature & SRTM \\
 & HAND (m) & MERIT Hydro \\
 & TWI, SPI & MERIT Hydro + SRTM \\
\midrule
Hydrological & Dist. to drainage, Drainage density & HydroSHEDS \\
\midrule
Climatic & Mean/Max precip., Soil moisture & CHIRPS, ERA5 \\
\midrule
Land/Demog. & Land cover, NDVI, Pop. density & WorldCover, S2, WorldPop \\
\midrule
Water hist. & JRC occurrence, SAR flood freq., Accessibility$^*$ & JRC, This study, Oxford MAP \\
\bottomrule
\multicolumn{3}{l}{\scriptsize $^*$Accessibility: travel time (min) to nearest city of $\geq$50,000 inhabitants \citep{Weiss2018}.}
\end{tabular}
\end{table}

\textbf{Training labels---adapted for arid regions.} Flood-positive pixels were defined as locations with JRC occurrence $\geq 10\%$ (lowered from the 25\% used in humid Antioquia \citep{EspinalMaya2026} to account for lower flood frequency) OR HAND $< \SI{5}{\metre}$. Flood-negative pixels required JRC $< 3\%$, HAND $> \SI{30}{\metre}$, and slope $> 10^{\circ}$, stratified across three subregions. The number of samples per class was set to 2,500 (reduced from 5,000) to match the smaller study area and lower flood prevalence. The balanced dataset was split 70/30 for training/test.

Three models were trained: Random Forest \citep{Breiman2001} (scikit-learn 1.5.2), XGBoost \citep{Chen2016} (XGBoost 2.1.3), and LightGBM \citep{Ke2017} (LightGBM 4.5.0). Hyperparameters were selected via randomized search (100 iterations, 3-fold inner CV on the training partition): Random Forest (500 trees, max\_depth $= 15$, min\_samples\_leaf $= 5$), XGBoost (500 estimators, learning\_rate $= 0.05$, max\_depth $= 8$, subsample $= 0.8$, colsample\_bytree $= 0.8$), and LightGBM (500 estimators, learning\_rate $= 0.05$, num\_leaves $= 63$, min\_child\_samples $= 20$). A weighted ensemble combined predictions proportional to each model's AUC-ROC (Equation~\ref{eq:ensemble}).

\begin{equation}
    P_{\text{ens}}(x) = \sum_{k=1}^{3} w_k \cdot P_k(x), \quad w_k = \frac{\text{AUC}_k}{\sum_{j} \text{AUC}_j}
    \label{eq:ensemble}
\end{equation}

Performance was evaluated via \textit{spatial} five-fold cross-validation \citep{Roberts2017}. With only three subregions, the 10 municipalities of Baja Guajira were divided into three geographic groups to yield five spatially contiguous folds: Fold~1 (Alta Guajira: Uribia, Manaure), Fold~2 (Media Guajira: Riohacha, Maicao, Dibulla), Fold~3 (Baja Guajira--North: Albania, Barrancas, Hatonuevo, Fonseca), Fold~4 (Baja Guajira--Central: Distracci\'on, El Molino, San Juan del Cesar), and Fold~5 (Baja Guajira--South: La Jagua del Pilar, Urumita, Villanueva). Metrics: AUC-ROC, accuracy, precision, recall, F1, Cohen's $\kappa$, Brier score. Feature importance was assessed using SHAP \citep{Lundberg2017}.

\subsection{Population Exposure and Risk Index}

Exposed population per municipality $j$ was computed as:
\begin{equation}
    \text{Pop}_{\text{exp},j} = \sum_{i \in \Omega_j} P_i \cdot \mathbb{1}[S_i \geq 0.6]
    \label{eq:pop_exposure}
\end{equation}
where $P_i$ is WorldPop population at pixel $i$ and $S_i$ is ensemble susceptibility. A composite Flood Risk Index integrates hazard, susceptibility, and exposure (Equation~\ref{eq:fri}):
\begin{equation}
    \text{FRI}_j = \frac{A_{\text{high},j}}{A_{\text{total},j}} \times \frac{\text{Pop}_{\text{exp},j}}{\text{Pop}_{\text{total},j}} \times FF_{\text{mean},j}
    \label{eq:fri}
\end{equation}
Sensitivity to the threshold $\tau$ was tested at 0.5, 0.6, and 0.7.

\subsection{Seasonal and ENSO Analysis}

Monthly water extent was stratified by ENSO phase using the Oceanic Ni\~no Index (ONI $> 0.5$\textdegree{}C for El Ni\~no, $< -0.5$\textdegree{}C for La Ni\~na, over $\geq$5 consecutive overlapping 3-month periods). ENSO years in the study period: El Ni\~no (2015, 2016, 2023, 2024), La Ni\~na (2020, 2021, 2022), Neutral (2017, 2018, 2019, 2025). Differences were tested via Kruskal-Wallis with Dunn's post-hoc correction. Precipitation trends were assessed via Mann-Kendall test \citep{Mann1945} with Sen's slope estimator \citep{Sen1968}. To characterize the drought--flood paradox, we computed antecedent soil moisture (ERA5-Land) at 1, 3, and 6-month lags and tested its correlation with flood extent.

%%%%%%%%%%%%%%%%%%%%%%%%%%%%%%%%%%%%%%%%%%
\section{Results}
\label{sec:results}

\subsection{SAR Water Detection and Sand Exclusion}

\begin{figure}[htbp]
\centering
\begin{minipage}[t]{0.48\textwidth}
\centering
\includegraphics[width=\textwidth]{fig03_jrc_water_occurrence.pdf}
\end{minipage}\hfill
\begin{minipage}[t]{0.48\textwidth}
\centering
\includegraphics[width=\textwidth]{fig05_hand_susceptibility.pdf}
\end{minipage}
\caption{(\textbf{Left})~JRC Global Surface Water occurrence for La Guajira, showing concentrated water dynamics along the R\'io Rancher\'ia corridor and Sierra Nevada piedmont rivers. (\textbf{Right})~HAND (Height Above Nearest Drainage) map; low HAND values (red) trace the fluvial network and indicate high flood susceptibility. Note the extensive low-HAND coastal plain in Alta Guajira, where the Sand Exclusion Layer prevents false flood classification.}
\label{fig:jrc_hand}
\end{figure}

\textit{Note: Quantitative results in this section are placeholders that will be updated upon running the GEE processing pipeline. The methodological framework and analytical structure are final.}

The Sand Exclusion Layer classified approximately [X]\% of the department's area as persistent sand (concentrated in Alta Guajira and the Manaure salt flats). Without the SEL, preliminary tests showed [X]\% false positive water detections in the sand-dominated areas---confirming the critical necessity of this correction for arid-zone SAR flood mapping.

Processing Sentinel-1 scenes yielded monthly composites and 11 annual maximum extent maps. Adaptive Otsu thresholds are expected to show stronger seasonal variation than in humid Antioquia, reflecting the sharper contrast between dry-season (near-zero surface water) and wet-season (concentrated flooding along the Rancher\'ia corridor) conditions.

The flood frequency analysis is expected to reveal a concentrated spatial pattern dominated by the R\'io Rancher\'ia corridor from Barrancas through Albania to the Riohacha delta, the Sierra Nevada piedmont rivers (Palomino, Jerez) in Dibulla, and ephemeral arroyos around Maicao and Uribia that activate only during intense rainfall events. Unlike Antioquia's distributed flooding across multiple river systems, La Guajira's flood-prone area is expected to be more spatially concentrated and temporally intermittent.

\subsection{Machine Learning Model Performance}

\begin{table}[htbp]
\caption{Flood susceptibility model performance under spatial five-fold cross-validation (mean $\pm$ SD). Results to be populated after pipeline execution.}
\label{tab:ml_results}
\centering
\footnotesize
\begin{tabular}{lccccccc}
\toprule
\textbf{Model} & \textbf{AUC-ROC} & \textbf{Accuracy} & \textbf{Precision} & \textbf{Recall} & \textbf{F1} & \textbf{$\kappa$} & \textbf{Brier} \\
\midrule
Random Forest   & -- & -- & -- & -- & -- & -- & -- \\
XGBoost         & -- & -- & -- & -- & -- & -- & -- \\
LightGBM        & -- & -- & -- & -- & -- & -- & -- \\
\textbf{Ensemble} & -- & -- & -- & -- & -- & -- & -- \\
\bottomrule
\end{tabular}
\end{table}

\begin{table}[htbp]
\caption{Confusion matrix for the ensemble model on the held-out test set. To be populated after pipeline execution.}
\label{tab:confusion}
\centering
\footnotesize
\begin{tabular}{lcc}
\toprule
 & \textbf{Predicted flood} & \textbf{Predicted non-flood} \\
\midrule
\textbf{Observed flood}     & -- (TP) & -- (FN) \\
\textbf{Observed non-flood} & -- (FP) & -- (TN) \\
\bottomrule
\end{tabular}
\end{table}

\begin{figure}[htbp]
\centering
\begin{minipage}[t]{0.46\textwidth}
\centering
\includegraphics[width=\textwidth]{fig06_roc_curves.pdf}
\end{minipage}\hfill
\begin{minipage}[t]{0.52\textwidth}
\centering
\includegraphics[width=\textwidth]{fig07_shap_importance.pdf}
\end{minipage}
\caption{(\textbf{Left})~ROC curves for individual models and the weighted ensemble under spatial cross-validation. (\textbf{Right})~Global SHAP feature importance (top 12 of 18 predictors). Based on the semi-arid literature \citep{Bammou2024, MoroccoSARML2024}, HAND, elevation, and distance to drainage are expected to dominate, with rainfall-related features potentially showing higher importance than in humid Antioquia due to the sharper precipitation gradient across subregions.}
\label{fig:roc_shap}
\end{figure}

Based on comparable semi-arid flood susceptibility studies \citep{Bammou2024, Ren2024, MoroccoSARML2024}, ensemble AUC-ROC values of 0.85--0.93 are expected. Two factors may produce slightly lower performance than the 0.94 achieved in Antioquia: (1) fewer flood events reduce the information content of SAR-derived features, and (2) the sand--water confusion, even after SEL correction, introduces residual noise in the training labels. Conversely, the sharper topographic contrast in Baja Guajira (Sierra Nevada foothills versus lowland desert) may enhance model discrimination in those folds.

\subsection{Feature Importance}

SHAP analysis is expected to reveal a hierarchy partially different from Antioquia. Based on the semi-arid ML literature:

\begin{itemize}
    \item \textbf{HAND} is expected to remain the dominant predictor, consistent with \citet{Bammou2024} and the Antioquia results (SHAP $= 0.182$). The sharp HAND gradient between Sierra Nevada piedmont and desert lowland should produce strong discrimination.
    \item \textbf{Elevation} may rank higher than in Antioquia because the study area spans from sea level to $>$\SI{5000}{\metre} (Sierra Nevada boundary) in a smaller area.
    \item \textbf{Annual precipitation} is expected to show higher SHAP importance than in Antioquia (where it ranked 11th), because the precipitation gradient across La Guajira (125--2,500+~mm\,yr$^{-1}$) is extreme and directly controls flood potential.
    \item \textbf{SAR flood frequency} may show lower importance due to the reduced number of flood events in the arid zones, limiting the feature's discriminative power in Alta Guajira.
    \item \textbf{Distance to drainage} should rank highly because flooding in La Guajira is tightly concentrated along defined channels rather than distributed across broad floodplains.
\end{itemize}

As in the Antioquia study \citep{EspinalMaya2026}, we will conduct an ablation experiment removing SAR flood frequency and JRC occurrence to assess circularity. Given the expected lower importance of these features in the arid context, the ablation-induced AUC reduction may be smaller than the 0.03 observed in Antioquia.

\subsection{Flood Susceptibility Map}

\begin{figure}[htbp]
\centering
\includegraphics[width=0.55\textwidth]{fig08_flood_susceptibility.pdf}
\caption{Ensemble flood susceptibility map for La Guajira. Dashed lines delineate the three subregions. The R\'io Rancher\'ia corridor and Sierra Nevada piedmont rivers are expected to concentrate the highest susceptibility values.}
\label{fig:susceptibility}
\end{figure}

The susceptibility map is expected to show a starkly different spatial pattern from Antioquia:

\begin{enumerate}
    \item \textbf{R\'io Rancher\'ia corridor} (Barrancas--Albania--Riohacha): The principal flood-susceptible zone, following HAND $< \SI{10}{\metre}$ terrain along the $\sim$\SI{248}{\kilo\metre} river course.
    \item \textbf{Sierra Nevada piedmont} (Dibulla): Flash-flood susceptibility along short, steep rivers (Palomino, Jerez, Ancho) descending from $>$\SI{3000}{\metre} to sea level in $<$\SI{50}{\kilo\metre}.
    \item \textbf{Riohacha--Manaure coastal plain}: Combined fluvial (Rancher\'ia delta), coastal (storm surge), and arroyo flooding.
    \item \textbf{Alta Guajira arroyos}: Localized susceptibility along ephemeral channels in Uribia and Manaure, activated during rare but intense rainfall events.
\end{enumerate}

A substantially larger proportion of the department is expected to fall in the very low/low susceptibility classes compared to Antioquia, reflecting the hyper-arid Alta Guajira where flooding is physically improbable except along defined arroyos.

\begin{table}[htbp]
\caption{Area distribution across susceptibility classes by subregion (\% of subregional area). To be populated after pipeline execution.}
\label{tab:suscept_area}
\centering
\footnotesize
\begin{tabular}{lrrrrr}
\toprule
\textbf{Subregion} & \textbf{Very low} & \textbf{Low} & \textbf{Moderate} & \textbf{High} & \textbf{Very high} \\
\midrule
Alta Guajira       & -- & -- & -- & -- & -- \\
Media Guajira      & -- & -- & -- & -- & -- \\
Baja Guajira       & -- & -- & -- & -- & -- \\
\midrule
\textbf{La Guajira} & -- & -- & -- & -- & -- \\
\bottomrule
\end{tabular}
\end{table}

\subsection{Population Exposure}

\begin{table}[htbp]
\caption{Population exposure by subregion ($\tau = 0.6$). FRI = Flood Risk Index. To be populated after pipeline execution.}
\label{tab:exposure}
\centering
\footnotesize
\begin{tabular}{lrrrrr}
\toprule
\textbf{Subregion} & \textbf{Area (km$^2$)} & \textbf{High susc. (\%)} & \textbf{Pop. ($\times 10^3$)} & \textbf{Pop. exp. ($\times 10^3$)} & \textbf{FRI} \\
\midrule
Alta Guajira       & $\sim$9,000  & -- & 270  & -- & -- \\
Media Guajira      & $\sim$5,000  & -- & 434  & -- & -- \\
Baja Guajira       & $\sim$7,000  & -- & 260  & -- & -- \\
\midrule
\textbf{Total}     & \textbf{20,848} & -- & \textbf{$\sim$964} & -- & -- \\
\bottomrule
\end{tabular}
\end{table}

\begin{figure}[htbp]
\centering
\begin{minipage}[t]{0.48\textwidth}
\centering
\includegraphics[width=\textwidth]{fig09_municipal_risk_ranking.pdf}
\end{minipage}\hfill
\begin{minipage}[t]{0.48\textwidth}
\centering
\includegraphics[width=\textwidth]{fig10_population_exposure.pdf}
\end{minipage}
\caption{(\textbf{Left})~All 15 municipalities ranked by Flood Risk Index; bar colour indicates subregion. (\textbf{Right})~Municipal-level population exposure to high/very high flood susceptibility zones; circle size proportional to exposed population, colour indicates proportion exposed.}
\label{fig:exposure_ranking}
\end{figure}

The municipalities with the highest expected FRI are Riohacha (Rancher\'ia delta, coastal flooding, population density), Dibulla (Sierra Nevada flash floods), and Barrancas (mid-Rancher\'ia flooding). A critical dimension of the exposure analysis is the Wayuu population: approximately 45\% of La Guajira's residents are Wayuu, concentrated in Alta and Media Guajira, where access to early warning systems, evacuation routes, and emergency services is extremely limited \citep{Contreras2020, FernandezLopera2024}. \citet{FernandezLopera2024} identified La Guajira as the Colombian department with the highest climate vulnerability index, driven by the intersection of physical hazard exposure and socioeconomic marginalization.

\subsection{Seasonal and ENSO Dynamics}

\begin{figure}[htbp]
\centering
\includegraphics[width=0.85\textwidth]{fig11_seasonal_dynamics.pdf}
\caption{Temporal flood dynamics (2015--2025). (\textbf{a})~Monthly flood extent with ENSO phase shading (blue: La Ni\~na; pink: El Ni\~no). (\textbf{b})~Climatological monthly mean showing unimodal pattern; September--November peak months highlighted.}
\label{fig:seasonal}
\end{figure}

Unlike Antioquia's bimodal seasonality, La Guajira's flood dynamics show a pronounced unimodal pattern with virtually zero flood extent during December--March (dry season) and a sharp September--November peak. CHIRPS analysis over 2015--2025 confirms the extreme ENSO sensitivity: El Ni\~no years averaged \SI{793}{mm/yr} while La Ni\~na years averaged \SI{904}{mm/yr} (a 14\% increase); the 2022 triple-dip La Ni\~na year recorded \SI{1196}{mm} and 2015 strong El Ni\~no only \SI{406}{mm} \citep{GallegoRevilla2022}. No statistically significant long-term precipitation trends were detected (Mann--Kendall $p > 0.08$ for all seasons), consistent with ENSO-dominated interannual variability. However, the land surface temperature shows a statistically significant declining trend of $-0.34$~\SI{}{\celsius/yr} (Mann--Kendall $\tau = -0.67$, $p = 0.005$), likely reflecting increased cloudiness during the La Ni\~na-dominated latter half of the study period.

\begin{table}[htbp]
\caption{Annual precipitation and flood extent statistics by ENSO phase (2015--2025). Flood extent values are preliminary and will be updated with SAR-derived results.}
\label{tab:enso}
\centering
\footnotesize
\begin{tabular}{lcccc}
\toprule
\textbf{ENSO phase} & \textbf{$n$ years} & \textbf{Mean precip.\ (mm/yr)} & \textbf{Std dev.\ (mm)} & \textbf{Key years} \\
\midrule
El Ni\~no  & 4 & 793 & 279 & 2015 (406~mm), 2016, 2023, 2024 \\
Neutral    & 4 & 702 & 101 & 2017, 2018, 2019, 2025 \\
La Ni\~na  & 3 & 904 & 285 & 2020, 2021, 2022 (1196~mm) \\
\bottomrule
\end{tabular}
\end{table}

\textbf{Drought--flood paradox.} A distinctive contribution of this study is the analysis of antecedent soil moisture conditions. During El Ni\~no droughts, soils in La Guajira become desiccated and surface-crusted, dramatically reducing infiltration capacity \citep{Alabbad2023}. We hypothesize that the transition from El Ni\~no to La Ni\~na produces disproportionately severe flooding because: (i) the rainfall increase is superimposed on soils with reduced permeability, and (ii) vegetation loss during drought exposes bare soil, increasing runoff coefficients. This mechanism---well-documented conceptually \citep{DroughtFloodMonitoring2024} but never quantified for La Guajira---may explain why the 2022 and 2024 flood events were so destructive despite La Guajira receiving less total rainfall than humid departments.

%%%%%%%%%%%%%%%%%%%%%%%%%%%%%%%%%%%%%%%%%%
\section{Discussion}
\label{sec:discussion}

\subsection{SAR Flood Mapping in Arid Environments: Methodological Contribution}

The Sand Exclusion Layer represents a critical methodological adaptation without which SAR-based flood mapping in La Guajira would produce unacceptable false positive rates. Our approach extends \citet{Martinis2018}, who developed the SEL concept for Somalia, and complements \citet{Garg2024}, who demonstrated coherence-amplitude fusion for arid flood detection. Three aspects merit discussion. First, the temporal persistence threshold (80\% of dry-season scenes) balances sand detection sensitivity with specificity---higher thresholds risk missing mobile dune fields, while lower thresholds may exclude legitimate flood-prone depressions. Second, the Manaure salt flat masking required special treatment: salt pans transition between dry (low backscatter = sand-like) and flooded (low backscatter = water-like) states on seasonal timescales, and the temporal statistics approach alone cannot distinguish these states. We resolved this by combining the SEL with JRC occurrence data and a static salt-pan mask derived from ESA WorldCover. Third, the minimum mapping unit was increased from \SI{1}{\hectare} (Antioquia) to \SI{2}{\hectare} to filter transient low-backscatter patches from small arroyo channels and wind-deposited sand features.

\subsection{Comparison with Antioquia Framework}

This study demonstrates the replicability of the framework developed by \citet{EspinalMaya2026} in a contrasting environment. Key differences are expected in four domains:

\begin{enumerate}
    \item \textbf{Spatial pattern:} Antioquia's flooding is distributed across multiple river systems (Cauca, Magdalena, Atrato, Nech\'i) with large alluvial floodplains. La Guajira's flooding is concentrated along a single dominant axis (R\'io Rancher\'ia) with flash-flood contributions from Sierra Nevada piedmont rivers.
    \item \textbf{Temporal dynamics:} Antioquia's bimodal seasonality (two distinct flood peaks per year) contrasts with La Guajira's unimodal pattern (single September--November peak), with near-zero flooding during the extended dry season.
    \item \textbf{SAR challenges:} In Antioquia, the primary SAR challenge is forest canopy interference (double-bounce scattering under dense vegetation). In La Guajira, it is the opposite: sand--water confusion in barren terrain. These complementary challenges test the framework's adaptability.
    \item \textbf{Feature importance:} Precipitation-related features are expected to rank substantially higher in La Guajira's SHAP analysis because rainfall is the limiting factor for flooding in arid environments, whereas in Antioquia, rainfall is abundant everywhere and topography is the primary discriminant.
\end{enumerate}

\subsection{The Drought--Flood Paradox}

La Guajira exemplifies a phenomenon well-documented in arid hydrology but rarely quantified with satellite data: the same region that experiences extreme drought also faces catastrophic flooding, and the drought itself exacerbates subsequent flood severity. The mechanism operates through three pathways: (1) \textit{soil crusting}---prolonged drying creates hydrophobic surface layers that repel initial rainfall, increasing surface runoff \citep{Alabbad2023}; (2) \textit{vegetation loss}---drought-induced plant mortality reduces canopy interception and root-zone infiltration; and (3) \textit{channel adjustment}---ephemeral arroyos accumulate sediment during dry periods, reducing channel capacity when flows return.

This paradox has direct policy implications. Traditional flood risk frameworks treat drought and flood as opposing hazards, but in La Guajira they are causally linked: investing in drought resilience (soil conservation, vegetation restoration, water harvesting) simultaneously reduces flood vulnerability by maintaining infiltration capacity and channel function. The conceptual framework of \citet{DroughtFloodMonitoring2024} for simultaneous drought--flood monitoring in Syria provides a methodological precedent, though our ENSO-stratified approach with antecedent moisture analysis extends it to the coupled ocean-atmosphere forcing that dominates La Guajira's climate.

\subsection{Wayuu Vulnerability and Environmental Justice}

The exposure analysis reveals a structural injustice: the Wayuu community, which has been identified as one of the most climate-vulnerable populations in Colombia \citep{Contreras2020, FernandezLopera2024, Contreras2019wayuu}, occupies territories in Alta and Media Guajira that are exposed to both drought and flood extremes. During the 2012--2016 drought, child mortality among the Wayuu reached 23.4 per 1,000---four times the national rate---driven by water scarcity, malnutrition, and limited healthcare access \citep{Contreras2020}. When La Ni\~na returns, these same communities face flash flooding along arroyos and coastal storm surge, often without access to early warning systems or evacuation infrastructure.

\citet{FernandezLopera2024} found that La Guajira has the highest climate vulnerability index of any Colombian department, and proposed culturally appropriate risk transfer mechanisms. Our municipality-scale flood susceptibility maps provide the missing physical hazard layer for such interventions: by identifying which Wayuu-inhabited territories face the highest compound exposure (drought $+$ flood), resources can be targeted to communities where resilience investments yield the greatest reduction in multi-hazard risk. The water governance framework proposed by Wayuu communities themselves \citep{WayuuWaterGovernance2025}---integrating traditional ecological knowledge with modern infrastructure---aligns with this approach.

\subsection{Limitations, Uncertainty, and Future Directions}
\label{sec:limitations}

Five limitations warrant quantification. (1) \textbf{Sand--water residual confusion:} Despite the SEL, mobile dune fields in Alta Guajira may produce transient low-backscatter signatures that evade the temporal persistence filter. Incorporating InSAR coherence \citep{Garg2024} would reduce this residual error but requires revisit-pair availability. (2) \textbf{Fewer flood events:} The reduced flood frequency in arid La Guajira means fewer positive training samples and potentially less stable model performance, particularly in Alta Guajira where some years may have zero detectable flooding. The reduction to 2,500 samples per class mitigates overfitting but may limit the model's ability to capture rare flash-flood patterns. (3) \textbf{SAR circularity:} As in the Antioquia study, SAR flood frequency appears as both a labelling criterion and a predictor. The same three safeguards apply (spatial CV, ablation experiment, conservative labelling threshold), though the expected lower importance of SAR-derived features in the arid context may further reduce circularity concerns. (4) \textbf{HAND accuracy in low-relief terrain:} Much of Media and Alta Guajira has terrain gradients approaching or below the SRTM vertical error ($\pm$\SI{5}{\metre}). In these flat areas, HAND-based susceptibility estimates carry greater uncertainty than in the mountainous Baja Guajira. (5) \textbf{Wayuu settlement mapping:} WorldPop \SI{100}{\metre} data may underestimate dispersed Wayuu rancheria settlement patterns, which are sparser and more mobile than assumed by standard population models. Integration with community-based mapping data \citep{Contreras2019wayuu} would improve exposure estimates.

Future priorities include: (1) InSAR coherence integration for improved arid flood detection; (2) coupling with the drought monitoring framework of \citet{MeraGomez2021} for a unified multi-hazard assessment; (3) high-resolution LiDAR DEM acquisition for the R\'io Rancher\'ia corridor; (4) near-real-time GEE flood monitoring dashboard linked to ENSO forecasts; and (5) participatory validation with Wayuu communities of flood susceptibility maps against traditional flood knowledge.

%%%%%%%%%%%%%%%%%%%%%%%%%%%%%%%%%%%%%%%%%%
\section{Conclusions}
\label{sec:conclusions}

This study demonstrates that the SAR-ensemble ML framework developed for humid tropical Antioquia \citep{EspinalMaya2026} can be successfully adapted to the contrasting arid/semi-arid environment of La Guajira---Colombia's most climate-vulnerable department---through targeted methodological modifications: a Sand Exclusion Layer for SAR, adjusted training thresholds for arid flood conditions, and coastal/tidal masking.

\textbf{Methodological contribution:} The Sand Exclusion Layer overcomes the fundamental limitation that has prevented SAR-based flood mapping in arid regions. By leveraging multi-temporal backscatter statistics, we resolve the sand--water confusion that causes global flood monitoring products to exclude arid areas entirely \citep{Garg2024}. The framework's replicability across two hydroclimatically contrasting departments---from humid Andean mountains to Caribbean desert---validates its generalizability across tropical environments.

\textbf{Scientific contribution:} We provide the first municipality-scale flood susceptibility map for a Colombian arid department, quantifying the drought--flood paradox through ENSO-stratified analysis. The framework demonstrates that antecedent drought conditions amplify subsequent flood severity---a mechanism with direct implications for multi-hazard early warning systems.

\textbf{Applied contribution:} The explicit assessment of Wayuu indigenous community exposure bridges the gap between physical hazard mapping and vulnerability-informed risk management. By producing municipality-scale outputs directly ingestible by Colombian POT and PMGRD instruments, the framework transforms satellite data into governance-actionable products in a department where such information has been entirely absent.

The paradox that defines La Guajira---extreme drought and catastrophic flooding within the same communities---demands integrated solutions. Flood risk cannot be managed in isolation from drought resilience, and neither can be managed without engaging the Wayuu communities who bear the disproportionate burden of both extremes. This framework provides the physical hazard layer; closing the governance gap requires pairing it with the culturally grounded adaptation strategies that communities have been developing for centuries \citep{WayuuWaterGovernance2025}.

%%%%%%%%%%%%%%%%%%%%%%%%%%%%%%%%%%%%%%%%%%
\section*{Supplementary Materials}

The following supplementary materials are available online: Table~S1 (Flood Risk Index for all 15 municipalities), Figures~S1--S3 (SHAP dependence plots for HAND, flood frequency, and elevation), Figure~S4 (per-fold ROC curves), Figure~S5 (Sand Exclusion Layer validation), and the complete Google Earth Engine processing scripts.

\section*{Author Contributions}

Conceptualization, C.E.M. and S.J.L.; methodology, C.E.M.; software, C.E.M.; formal analysis, C.E.M.; investigation, C.E.M.; data curation, C.E.M.; writing---original draft, C.E.M.; writing---review and editing, S.J.L.; visualization, C.E.M.; supervision, S.J.L. All authors have read and agreed to the published version of the manuscript.

\section*{Funding}

This research was supported by Universidad EAFIT through the Master's thesis research program. No external funding was received.

\section*{Data Availability}

All satellite data were accessed through Google Earth Engine. Administrative boundaries from FAO GAUL 2015. Source code available at \url{https://github.com/Cespial/guajira-flood-risk} (MIT license). The Antioquia framework on which this study is based is available at \url{https://github.com/Cespial/antioquia-flood-risk}.

\section*{Acknowledgments}

The authors thank Google Earth Engine for cloud computing access, ESA for Sentinel-1 data, the JRC for Global Surface Water, UNGRD for maintaining Colombia's disaster event database, and the Wayuu communities of La Guajira whose resilience in the face of compound climate extremes motivates this work.

\section*{Conflicts of Interest}

The authors declare no conflicts of interest.

\section*{Abbreviations}

The following abbreviations are used in this manuscript:\\

\noindent
\begin{tabular}{@{}ll}
SAR & Synthetic Aperture Radar\\
GEE & Google Earth Engine\\
HAND & Height Above Nearest Drainage\\
ENSO & El Ni\~no--Southern Oscillation\\
SHAP & SHapley Additive exPlanations\\
FRI & Flood Risk Index\\
SEL & Sand Exclusion Layer\\
POT & Plan de Ordenamiento Territorial\\
PMGRD & Plan Municipal de Gesti\'on del Riesgo de Desastres\\
AUC-ROC & Area Under the ROC Curve\\
TWI & Topographic Wetness Index\\
SPI & Stream Power Index\\
ONI & Oceanic Ni\~no Index
\end{tabular}

%%%%%%%%%%%%%%%%%%%%%%%%%%%%%%%%%%%%%%%%%%

\bibliographystyle{unsrtnat}
\bibliography{references}

\end{document}
